\documentclass[10pt,a4paper]{amsart}
\usepackage[margin=19mm]{geometry}
\usepackage{amsmath,amssymb,mathtools}
\usepackage[scr=rsfs,cal=euler]{mathalfa}
\usepackage{xfrac}
\usepackage[unicode,hidelinks,bookmarks=false]{hyperref}
\newcommand{\ie}{i.e.\ }
\newcommand{\cf}{cf.\ }
\newcommand{\resp}{respectively\ }
\newcommand{\Subsection}{\subsection}
\providecommand{\nobreakdash}{\nobreak\hbox{-}}
\setlength{\emergencystretch}{2em}
\multlinegap0pt
\usepackage{bm}
\usepackage{bbm}
\usepackage{dsfont}
\usepackage{xspace}
\usepackage{cancel}
\usepackage{mathtools}
\usepackage{amsthm}
\makeatletter
\@ifundefined{lemma}{\newtheorem{lemma}{Lemma}}{}
\makeatother
\makeatletter
\expandafter\ifx\csname claim\endcsname\relax
\newenvironment{claim}{{\em Claim:}\/}{}
\fi
\expandafter\ifx\csname proofclaim\endcsname\relax
\newenvironment{proofclaim}{{\em Proof:}\/}{}
\fi
\makeatother
\title[Modal definability in Euclidean modal logics]{Modal definability in Euclidean modal logics}
\author{Philippe Balbiani, Tinko Tinchev}
\date{Source: arXiv:2508.10813v1 (CC BY 4.0)}
\begin{document}
\maketitle

\noindent\emph{Source and licence.} Modal definability in Euclidean modal logics. Version arXiv:2508.10813v1 (CC BY 4.0), distributed under the Creative Commons Attribution 4.0 International licence, as recorded in the arXiv licence metadata for this exact version and shown on the abstract page https://arxiv.org/abs/2508.10813v1. Full rights notice: licenses/arxiv-2508.10813-CC-BY-4.0.txt.\par\smallskip

\noindent\textit{Editorial scope.} Counted unit: the Lemma whose source label is \texttt{lemma:relatively:elementary:definable:K2} in the article (source0.tex lines 6680--6686, source label \texttt{lemma:relatively:elementary:definable:K2}), delivered together with the article's own complete original proof of it (source0.tex lines 6689--7280, one \texttt{proof} environment of 592 lines). No mathematics is written, repaired or re-derived: statement and proof are the article's own text line for line. Required context retained uncounted: none. Every definition, display and label of those lines is retained unchanged. Omitted from this package and not redistributed: 1--6679, 6687--6688, 7281--10629. Nothing inside a retained excerpt is omitted or altered: every retained body is delivered exactly as the article's lines are written, so no omission receipt of any kind accompanies this package. Citations: the retained excerpts carry no citation command at all, so no bibliography entry of the article is transcribed and no reference is redistributed. Reference adaptation: none was needed and none was made; every reference inside the delivered unit resolves inside it, so no unresolved reference or citation is produced. Printed slips kept literal: no printed word, symbol, hypothesis or number of the counted statement or of its proof was altered. Layout and class: the article's own class is \texttt{article} with packages including amsfonts, amsmath, amssymb, color, graphicx, latexsym; the wrapper is the lane's pinned \texttt{amsart} editorial wrapper at 10pt on A4, which restates the theorem-like environments the retained text needs and adds the article's own macro definitions that the retained text needs (none), with extra wrapper packages (bm, bbm, dsfont, xspace, cancel, mathtools). The delivered numbering is the unit's own. Assets: no retained span contains a figure environment, a graphics-inclusion command, a PDF-page inclusion, a table or a TikZ picture; no image file is copied into this package, the package declares no asset dependency and no image file is redistributed with it. Licence: version arXiv:2508.10813v1 of this article is distributed under the Creative Commons Attribution 4.0 International licence, as recorded in the arXiv licence metadata for this exact version and shown on the abstract page https://arxiv.org/abs/2508.10813v1. A full rights notice is delivered at licenses/arxiv-2508.10813-CC-BY-4.0.txt. Characters: every retained excerpt is pure ASCII, so the delivered engine, XeLaTeX with its default Latin Modern text font, reports no missing character.
\begin{lemma}\label{lemma:relatively:elementary:definable:K2}
%
%
The class of all irreflexive symmetric frames with at least $2$~elements is relatively elementary definable in ${\mathcal K}_{2}$.
%
%
\end{lemma}

\begin{proof}
%
%
Let $(W,R)$ be an arbitrary irreflexive symmetric frame with at least $2$ elements.
Let $E$ be $\{\{s_{1},s_{2}\}:\ s_{1},s_{2}{\in}W$ and $s_{1}{R}s_{2}\}$.
Notice that for all $s_{1},s_{2}{\in}W$, if $\{s_{1},s_{2}\}{\in}E$ then $s_{1}{\not=}s_{2}$.
Without loss of generality, we can assume that $E$ and $W{\times}\{(3,4)\}$ are disjoint.
Let $A$ be $E{\cup}(W{\times}\{3,4\})$.
Let $B$ be $W{\times}\{1,2\}$.
Let ${\mathcal F}^{\rho}_{A,B}$ be the galaxy such that
%
%
\begin{itemize}
%
%
\item for all $\{s_{1},s_{2}\}{\in}E$, $\rho(\{s_{1},s_{2}\}){=}\{(s_{1},1),(s_{2},1)\}$,
%
%
\item for all $s{\in}W$, $\rho((s,3)){=}\{(s,1),(s,2)\}$,
%
%
\item for all $s{\in}W$, $\rho((s,4)){=}\{(s,1),(s,2)\}$.
%
%
\end{itemize}
%
%
%Since $\{2\}{\times}\mathbf{N}^{-}{\subseteq}\mathtt{S}_{\L}$, then ${\mathcal F}^{\rho}_{A,B}{\models}\L$.
Obviously, ${\mathcal F}^{\rho}_{A,B}$ is in ${\mathcal K}_{2}$.
%
%
\\
\\
%
%
Let $X$ be $\{(a_{1},a_{2}):\ a_{1},a_{2}{\in}W^{\rho}_{A,B}$ and ${\mathcal F}^{\rho}_{A,B}{\models}\mathbf{U}(\mathbf{x_{1}},\mathbf{x_{2}})\ \lbrack a_{1},a_{2}\rbrack\}$ where
%
%
\begin{itemize}
%
%
\item $\mathbf{U}(\mathbf{x_{1}},\mathbf{x_{2}})::=\mathbf{x_{1}}{\not=}\mathbf{x_{2}}\wedge\mathbf{R}(\mathbf{x_{1}},\mathbf{x_{2}})\wedge\mathbf{R}(\mathbf{x_{2}},\mathbf{x_{1}})\wedge\exists^{{=}2}\mathbf{y}(\mathbf{R}(\mathbf{y}){=}\{\mathbf{x_{1}},\mathbf{x_{2}}\}\wedge\neg\exists\mathbf{z}\mathbf{R}(\mathbf{z},\mathbf{y}))$.
%
%
\end{itemize}
%
%
Obviously, for all $s{\in}W$, $((s,1),(s,2)),((s,2),(s,1)){\in}X$.
Hence, $X$ is non-empty.
%
%
\\
\\
%
%
Obviously, for all $(a_{1},a_{2}){\in}X$, there exists $s{\in}W$ such that $\{a_{1},a_{2}\}{=}\{(s,1),(s,
%
%
%
%
%
%
%
%
%
%
%
%
$\linebreak$
%
%
%
%
%
%
%
%
%
%
%
%
2)\}$.
%
%
\\
\\
%
%
Let $\bowtie$ be the binary relation on $X$ such that for all $(a_{1},a_{2}),(b_{1},b_{2}){\in}X$, $(a_{1},a_{2}){\bowtie}
%
%
%
%
%
%
%
%
%
%
%
%
$\linebreak$
%
%
%
%
%
%
%
%
%
%
%
%
(b_{1},b_{2})$ if and only if ${\mathcal F}^{\rho}_{A,B}{\models}\mathbf{E}(\mathbf{x_{1}},\mathbf{x_{2}},\mathbf{y_{1}},\mathbf{y_{2}})\ \lbrack a_{1},a_{2},b_{1},b_{2}\rbrack$ where
%
%
\begin{itemize}
%
%
\item $\mathbf{E}(\mathbf{x_{1}},\mathbf{x_{2}},\mathbf{y_{1}},\mathbf{y_{2}})::=\mathbf{U}(\mathbf{x_{1}},\mathbf{x_{2}})\wedge\mathbf{U}(\mathbf{y_{1}},\mathbf{y_{2}})\wedge\{\mathbf{x_{1}},\mathbf{x_{2}}\}{=}\{\mathbf{y_{1}},\mathbf{y_{2}}\}$.
%
%
\end{itemize}
%
%
Obviously, $\bowtie$ is an equivalence relation on $X$.
For all $(a_{1},a_{2}){\in}X$, the equivalence class of $(a_{1},a_{2})$ modulo $\bowtie$ will be denoted $\lbrack(a_{1},a_{2})\rbrack$.
%
%
\\
\\
%
%
Since for all $(a_{1},a_{2}){\in}X$, there exists $s{\in}W$ such that $\{a_{1},a_{2}\}{=}\{(s,1),(s,2)\}$, then for all $(a_{1},a_{2}),(b_{1},b_{2}){\in}X$, $(a_{1},a_{2}){\bowtie}(b_{1},b_{2})$ if and only if there exists $s{\in}W$ such that $\{a_{1},a_{2}\}{=}\{(s,1),(s,2)\}$ and $\{b_{1},b_{2}\}{=}\{(s,1),(s,2)\}$.
%
%
\\
\\
%
%
Let $S$ be the binary relation on $X$ such that for all $(a_{1},a_{2}),(b_{1},b_{2}){\in}X$, $(a_{1},a_{2}){S}
%
%
%
%
%
%
%
%
%
%
%
%
$\linebreak$
%
%
%
%
%
%
%
%
%
%
%
%
(b_{1},b_{2})$ if and only if ${\mathcal F}^{\rho}_{A,B}{\models}\mathbf{A}(\mathbf{x_{1}},\mathbf{x_{2}},\mathbf{y_{1}},\mathbf{y_{2}})\ \lbrack a_{1},a_{2},b_{1},b_{2}\rbrack$ where
%
%
\begin{itemize}
%
%
\item $\mathbf{A}(\mathbf{x_{1}},\mathbf{x_{2}},\mathbf{y_{1}},\mathbf{y_{2}})::=\mathbf{U}(\mathbf{x_{1}},\mathbf{x_{2}})\wedge\mathbf{U}(\mathbf{y_{1}},\mathbf{y_{2}})\wedge\{\mathbf{x_{1}},\mathbf{x_{2}}\}{\cap}\{\mathbf{y_{1}},\mathbf{y_{2}}\}{=}\emptyset\wedge
%
%
%
%
%
%
%
%
%
%
%
%
$\linebreak$
%
%
%
%
%
%
%
%
%
%
%
%
\bigvee\{\exists^{{=}1}\mathbf{z}(\mathbf{R}(\mathbf{z}){=}\{\mathbf{x_{i}},\mathbf{y_{j}}\}\wedge\neg\exists\mathbf{t}\mathbf{R}(\mathbf{t},\mathbf{z})):\ 1{\leq}i,j{\leq}2\}$.
%
%
\end{itemize}
%
%
We claim that for all $(a_{1},a_{2}),(b_{1},b_{2}),(c_{1},c_{2}),(d_{1},d_{2}){\in}X$, if $(a_{1},a_{2}){\bowtie}(c_{1},c_{2})$ and $(b_{1},b_{2}){\bowtie}(d_{1},d_{2})$ then $(a_{1},a_{2}){S}(b_{1},b_{2})$ if and only if $(c_{1},c_{2}){S}(d_{1},d_{2})$.
If not, let $(a_{1},a_{2}),(b_{1},b_{2}),(c_{1},c_{2}),(d_{1},d_{2}){\in}X$ be such that $(a_{1},a_{2}){\bowtie}(c_{1},c_{2})$, $(b_{1},b_{2}){\bowtie}(d_{1},
%
%
%
%
%
%
%
%
%
%
%
%
$\linebreak$
%
%
%
%
%
%
%
%
%
%
%
%
d_{2})$ and either $(a_{1},a_{2}){S}(b_{1},b_{2})$ and not $(c_{1},c_{2}){S}(d_{1},d_{2})$, or not $(a_{1},a_{2}){S}(b_{1},b_{2})$ and $(c_{1},c_{2}){S}(d_{1},d_{2})$.
Without loss of generality, suppose $(a_{1},a_{2}){S}(b_{1},b_{2})$ and not $(c_{1},c_{2}){S}(d_{1},d_{2})$.
Thus, ${\mathcal F}^{\rho}_{A,B}{\models}\mathbf{A}(\mathbf{x_{1}},\mathbf{x_{2}},\mathbf{y_{1}},\mathbf{y_{2}})\ \lbrack a_{1},a_{2},b_{1},b_{2}\rbrack$ and ${\mathcal F}^{\rho}_{A,B}{\not\models}
%
%
%
%
%
%
%
%
%
%
%
%
$\linebreak$
%
%
%
%
%
%
%
%
%
%
%
%
\mathbf{A}(\mathbf{x_{1}},\mathbf{x_{2}},\mathbf{y_{1}},\mathbf{y_{2}})\ \lbrack c_{1},c_{2},d_{1},d_{2}\rbrack$.
Consequently,
%
%
\begin{description}
%
%
\item[$(C_{1})$] $\{a_{1},a_{2}\}{\cap}\{b_{1},b_{2}\}{=}\emptyset$,
%
%
\item[$(C_{2})$] there exists $i,j{\in}\{1,2\}$ such that for exactly $1$~$e$ in ${\mathcal F}_{A,B}^{\rho}$, $R_{A,B}^{\rho}(e){=}\{a_{i},b_{j}\}$ and ${R_{A,B}^{\rho}}^{{-}1}(e){=}\emptyset$,
%
%
\item[$(C_{3})$] at least one of the following conditions holds:
%
%
\begin{description}
%
%
\item[$(C_{3}^{a})$] $\{c_{1},c_{2}\}{\cap}\{d_{1},d_{2}\}{\not=}\emptyset$,
%
%
\item[$(C_{3}^{b})$] for all $k,l{\in}\{1,2\}$, either for all $f$ in ${\mathcal F}_{A,B}^{\rho}$, if $R_{A,B}^{\rho}(f){=}\{c_{k},d_{l}\}$ then ${R_{A,B}^{\rho}}^{{-}1}(f){\not=}\emptyset$, or there exists $f^{\prime},f^{\prime\prime}$ in ${\mathcal F}_{A,B}^{\rho}$ such that $f^{\prime}{\not=}f^{\prime\prime}$,
%
%
%
%
%
%
%
%
%
%
%
%
\linebreak$
%
%
%
%
%
%
%
%
%
%
%
%
R_{A,B}^{\rho}(f^{\prime}){=}\{c_{k},d_{l}\}$, $R_{A,B}^{\rho}(f^{\prime\prime}){=}\{c_{k},d_{l}\}$, ${R_{A,B}^{\rho}}^{{-}1}(f^{\prime}){=}\emptyset$ and
%
%
%
%
%
%
%
%
%
%
%
%
\linebreak$
%
%
%
%
%
%
%
%
%
%
%
%
{R_{A,B}^{\rho}}^{{-}1}(f^{\prime\prime}){=}\emptyset$.
%
%
\end{description}
%
%
\end{description}
%
%
%Since $(a_{1},a_{2}),(b_{1},b_{2}),(c_{1},c_{2}),(d_{1},d_{2}){\in}X$, ${\mathcal F}^{\rho}_{A,B}{\models}\mathbf{U}(\mathbf{x_{1}},\mathbf{x_{2}})\ \lbrack a_{1},a_{2}\rbrack$, ${\mathcal F}^{\rho}_{A,B}{\models}\mathbf{U}(\mathbf{x_{1}},\mathbf{x_{2}})\ \lbrack b_{1},b_{2}\rbrack$, ${\mathcal F}^{\rho}_{A,B}{\models}\mathbf{U}(\mathbf{x_{1}},\mathbf{x_{2}})\ \lbrack c_{1},c_{2}\rbrack$ and ${\mathcal F}^{\rho}_{A,B}{\models}\mathbf{U}(\mathbf{x_{1}},\mathbf{x_{2}})\ \lbrack d_{1},d_{2}\rbrack$.
Let $s,t,u,v{\in}W$ be such that $\{a_{1},a_{2}\}{=}\{(s,1),(s,2)\}$, $\{b_{1},b_{2}\}{=}\{(t,1),(t,2)\}$, $\{c_{1},c_{2}\}{=}\{(u,1),(u,2)\}$ and $\{d_{1},d_{2}\}{=}\{(v,1),(v,2)\}$.
Since $(a_{1},a_{2}){\bowtie}(c_{1},c_{2})$ and $(b_{1},b_{2}){\bowtie}(d_{1},d_{2})$, then ${\mathcal F}^{\rho}_{A,B}{\models}\mathbf{E}(\mathbf{x_{1}},\mathbf{x_{2}},\mathbf{y_{1}},\mathbf{y_{2}})\ \lbrack a_{1},a_{2},c_{1},c_{2}\rbrack$ and ${\mathcal F}^{\rho}_{A,B}{\models}\mathbf{E}(\mathbf{x_{1}},
%
%
%
%
%
%
%
%
%
%
%
%
$\linebreak$
%
%
%
%
%
%
%
%
%
%
%
%
\mathbf{x_{2}},\mathbf{y_{1}},\mathbf{y_{2}})\ \lbrack b_{1},b_{2},d_{1},d_{2}\rbrack$.
Hence, %${\mathcal F}^{\rho}_{A,B}{\models}\mathbf{U}(\mathbf{x_{1}},\mathbf{x_{2}})\ \lbrack a_{1},a_{2}\rbrack$, ${\mathcal F}^{\rho}_{A,B}{\models}\mathbf{U}(\mathbf{x_{1}},\mathbf{x_{2}})\ \lbrack b_{1},b_{2}\rbrack$, ${\mathcal F}^{\rho}_{A,B}{\models}\mathbf{U}(\mathbf{x_{1}},\mathbf{x_{2}})\ \lbrack c_{1},c_{2}\rbrack$ and ${\mathcal F}^{\rho}_{A,B}{\models}\mathbf{U}(\mathbf{x_{1}},\mathbf{x_{2}})\ \lbrack d_{1},d_{2}\rbrack$.
%Moreover,
$\{a_{1},a_{2}\}{=}\{c_{1},c_{2}\}$ and $\{b_{1},b_{2}\}{=}\{d_{1},d_{2}\}$.
Since $\{a_{1},a_{2}\}{=}\{(s,1),(s,2)\}$, $\{b_{1},b_{2}\}{=}\{(t,1),(t,2)\}$, $\{c_{1},c_{2}\}{=}\{(u,1),(u,2)\}$ and $\{d_{1},
%
%
%
%
%
%
%
%
%
%
%
%
$\linebreak$
%
%
%
%
%
%
%
%
%
%
%
%
d_{2}\}{=}\{(v,1),(v,2)\}$, then $s{=}u$ and $t{=}v$.
Moreover, by $(C_{1})$, $s{\not=}t$.
In other respect, by $(C_{2})$, there exists $k^{\prime},l^{\prime}{\in}\{1,2\}$ such that $a_{i}{=}(s,k^{\prime})$, $b_{j}{=}(t,l^{\prime})$ and for exactly $1$~$e$ in ${\mathcal F}_{A,B}^{\rho}$, $R_{A,B}^{\rho}(e){=}\{(s,k^{\prime}),(t,l^{\prime})\}$ and ${R_{A,B}^{\rho}}^{{-}1}(e){=}\emptyset$.
Since $\{c_{1},c_{2}\}{=}\{(u,1),(u,2)\}$, $s{=}u$, $\{d_{1},d_{2}\}{=}\{(v,1),(v,2)\}$, $t{=}v$ and $s{\not=}t$, then $\{c_{1},
%
%
%
%
%
%
%
%
%
%
%
%
$\linebreak$
%
%
%
%
%
%
%
%
%
%
%
%
c_{2}\}{\cap}\{d_{1},d_{2}\}{=}\emptyset$.
Moreover, there exists $k,l{\in}\{1,2\}$ such that $c_{k}{=}(s,k^{\prime})$ and $d_{l}{=}(t,l^{\prime})$.
Hence, $(C_{3}^{a})$ does not hold.
Thus, by $(C_{3})$, $(C_{3}^{b})$ holds and either for all $f$ in ${\mathcal F}_{A,B}^{\rho}$, if $R_{A,B}^{\rho}(f){=}\{(s,k^{\prime}),(t,l^{\prime})\}$ then ${R_{A,B}^{\rho}}^{{-}1}(f){\not=}\emptyset$, or there exists $f^{\prime},f^{\prime\prime}$ in ${\mathcal F}_{A,B}^{\rho}$ such that $f^{\prime}{\not=}f^{\prime\prime}$, $R_{A,B}^{\rho}(f^{\prime}){=}\{(s,k^{\prime}),(t,l^{\prime})\}$, $R_{A,B}^{\rho}(f^{\prime\prime}){=}\{(s,k^{\prime}),(t,l^{\prime})\}$, ${R_{A,B}^{\rho}}^{{-}1}(f^{\prime}){=}\emptyset$ and ${R_{A,B}^{\rho}}^{{-}1}(f^{\prime\prime}){=}\emptyset$.
In the former case, since $R_{A,B}^{\rho}(e){=}\{(s,k^{\prime}),
%
%
%
%
%
%
%
%
%
%
%
%
$\linebreak$
%
%
%
%
%
%
%
%
%
%
%
%
(t,l^{\prime})\}$, then ${R_{A,B}^{\rho}}^{{-}1}(e){\not=}\emptyset$: a contradiction.
In the latter case, since $a_{i}{=}(s,k^{\prime})$, $b_{j}{=}(t,l^{\prime})$, then there exists $f^{\prime},f^{\prime\prime}$ in ${\mathcal F}_{A,B}^{\rho}$ such that $f^{\prime}{\not=}f^{\prime\prime}$, $R_{A,B}^{\rho}(f^{\prime}){=}\{a_{i},b_{j}\}$, $R_{A,B}^{\rho}(f^{\prime\prime}){=}\{a_{i},b_{j}\}$, ${R_{A,B}^{\rho}}^{{-}1}(f^{\prime}){=}\emptyset$ and ${R_{A,B}^{\rho}}^{{-}1}(f^{\prime\prime}){=}\emptyset$: a contradiction.
Consequently, for all $(a_{1},a_{2}),(b_{1},b_{2}),(c_{1},c_{2}),(d_{1},d_{2}){\in}X$, if $(a_{1},a_{2}){\bowtie}(c_{1},c_{2})$ and $(b_{1},b_{2}){\bowtie}(d_{1},d_{2})$ then $(a_{1},a_{2}){S}(b_{1},b_{2})$ if and only if $(c_{1},c_{2}){S}(d_{1},d_{2})$.
%
%
\\
\\
%
%
Let $(W^{\prime},R^{\prime})$ be the quotient of $(X,S)$ modulo $\bowtie$~---~i.e.
%
%
\begin{itemize}
%
%
\item $W^{\prime}$ is the quotient set of $X$ modulo $\bowtie$,
%
%
\item $R^{\prime}$ is the binary relation on $W^{\prime}$ such that for all $(a_{1},a_{2}),(b_{1},b_{2}){\in}X$, $\lbrack(a_{1},a_{2})\rbrack{R^{\prime}}\lbrack(b_{1},b_{2})\rbrack$ if and only if $(a_{1},a_{2}){S}(b_{1},b_{2})$.
%
%
\end{itemize}
%
%
Since for all $(a_{1},a_{2}){\in}X$, there exists $s{\in}W$ such that $\{a_{1},a_{2}\}{=}\{(s,1),(s,2)\}$ and for all $(a_{1},a_{2}),(b_{1},b_{2}){\in}X$, $(a_{1},a_{2}){\bowtie}(b_{1},b_{2})$ if and only if there exists $s{\in}W$ such that $\{a_{1},a_{2}\}{=}\{(s,1),(s,2)\}$ and $\{b_{1},b_{2}\}{=}\{(s,1),(s,2)\}$, then let $\kappa:\ W^{\prime}{\longrightarrow}W$ be the function such that for all $(a_{1},a_{2}){\in}X$, $\kappa(\lbrack(a_{1},a_{2})\rbrack)$ is the unique $s{\in}W$ such that $\{a_{1},a_{2}\}{=}\{(s,1),(s,2)\}$.
Obviously, $\kappa$ is bijective.
%for all $s{\in}W$, $((s,1),(s,2)),((s,2),(s,1)){\in}X$ and
%
\\
\\
%
%
We claim that for all $(a_{1},a_{2}),(b_{1},b_{2}){\in}X$, $\lbrack(a_{1},a_{2})\rbrack{R^{\prime}}\lbrack(b_{1},b_{2})\rbrack$ if and only if $\kappa(\lbrack(a_{1},a_{2})\rbrack){R}\kappa(\lbrack(b_{1},b_{2})\rbrack)$.
Let $(a_{1},a_{2}),(b_{1},b_{2}){\in}X$.
Let $s$ be the unique element in $W$ such that $\{a_{1},a_{2}\}{=}\{(s,1),(s,2)\}$ and $t$ be the unique element in $W$ such that $\{b_{1},b_{2}\}{=}\{(t,1),(t,2)\}$.
Obviously, the following conditions are equivalent:
%
%
\begin{itemize}
%
%
\item $\lbrack(a_{1},a_{2})\rbrack{R^{\prime}}\lbrack(b_{1},b_{2})\rbrack$,
%
%
\item $(a_{1},a_{2}){S}(b_{1},b_{2})$,
%
%
\item ${\mathcal F}^{\rho}_{A,B}{\models}\mathbf{A}(\mathbf{x_{1}},\mathbf{x_{2}},\mathbf{y_{1}},\mathbf{y_{2}})\ \lbrack a_{1},a_{2},b_{1},b_{2}\rbrack$,
%
%
\item $s{\not=}t$ and there exists $i,j{\in}\{1,2\}$ such that for exactly $1$~$e$ in ${\mathcal F}_{A,B}^{\rho}$, $R_{A,B}^{\rho}(e){=}
%
%
%
%
%
%
%
%
%
%
%
%
$\linebreak$
%
%
%
%
%
%
%
%
%
%
%
%
\{a_{i},b_{j}\}$ and ${R_{A,B}^{\rho}}^{{-}1}(e){=}\emptyset$,
%
%
\item $\kappa(\lbrack(a_{1},a_{2})\rbrack){R}\kappa(\lbrack(b_{1},b_{2})\rbrack)$.
%
%
\end{itemize}
%
%
Hence, for all $(a_{1},a_{2}),(b_{1},b_{2}){\in}X$, $\lbrack(a_{1},a_{2})\rbrack{R^{\prime}}\lbrack(b_{1},b_{2})\rbrack$ if and only if $\kappa(\lbrack(a_{1},a_{2})\rbrack){R}
%
%
%
%
%
%
%
%
%
%
%
%
$\linebreak$
%
%
%
%
%
%
%
%
%
%
%
%
\kappa(\lbrack(b_{1},b_{2})\rbrack)$.
%
%
\\
\\
%
%
All in all, we have proved that, $(W^{\prime},R^{\prime})$ is isomorphic with $(W,R)$.
%
%
\\
\\
%
%
Thus, we have shown that the class of all irreflexive symmetric frames with at least $2$ elements is relatively elementary definable in ${\mathcal K}_{2}$.
%
%
\medskip
%
%
\end{proof}

\end{document}
